\section{Natural units}\label{app:A.1}

We set $\hbar = c = k_{\mathrm{B}} = 1$. Energy, mass, momentum, frequency and temperature are then all measured in electronvolts, length and time in $\mathrm{eV}^{-1}$, and speeds are pure numbers, fractions of the speed of light. A dark-matter speed of \SI{250.5}{km.s^{-1}} is $8.356\times10^{-4}$, and the energy and momentum it carries at mass $m_\chi$ are $\tfrac12 m_\chi v^2$ and $m_\chi v$ in the same unit. The choice suits a subject in which a particle-physics mass, a crystal momentum, a phonon frequency and a temperature enter one energy-conservation condition. Particle physics uses it routinely, and Ref.~\cite[Appendix~A, Table~III]{lin2019tasi} tabulates the astrophysical units it is mixed with; condensed-matter texts often keep $\hbar$ and $k_{\mathrm{B}}$, and their formulas are converted by restoring the factors below.

\subsection{Restoring the constants}

A formula in natural units is converted to conventional units by inserting the powers of $\hbar$, $c$ and $k_{\mathrm{B}}$ that make every term dimensionally consistent. Four constants suffice; their values follow from the 2022 adjustment of CODATA (the Committee on Data of the International Science Council)~\cite{mohr2025codata}, in which $c$, $h$, $e$ and $k_{\mathrm{B}}$ are exact:
\begin{equation}
  \begin{aligned}
  \hbar c &= \SI{1.973270e-5}{eV.cm} = \SI{197.3270}{MeV.fm}, &\quad \hbar &= \SI{6.582120e-16}{eV.s},\\
  k_{\mathrm{B}} &= \SI{8.617333e-5}{eV.K^{-1}}, &\quad c &= \SI{2.997925e5}{km.s^{-1}}
  \end{aligned}.
  \label{eq:A.1-constants}
\end{equation}
A length $L$ becomes $L/(\hbar c)$ in $\mathrm{eV}^{-1}$, a time $t$ becomes $t/\hbar$, a temperature $T$ becomes $k_{\mathrm{B}}T$ and a speed $u$ becomes $u/c$. A density per unit volume picks up $(\hbar c)^3$, a cross-section $(\hbar c)^{-2}$, and a rate per unit time $\hbar$.

\begin{table}[t]
  \centering
  \caption{Conversion factors between conventional units and natural units ($\hbar = c = k_{\mathrm{B}} = 1$), computed from CODATA 2022 values~\cite{mohr2025codata}. Examples in the last column are values used in the book.}
  \label{tab:A-conversions}
  \small
  \begin{tabular}{@{}p{0.17\linewidth}p{0.42\linewidth}p{0.35\linewidth}@{}}
    \toprule
    Quantity & Conversion & Example \\
    \midrule
    length & $\SI{1}{cm} = \SI{5.0677e4}{eV^{-1}}$, $\SI{1}{\angstrom} = \SI{5.0677e-4}{eV^{-1}}$ & Bohr radius $1/(\alpha m_e) = \SI{0.5292}{\angstrom}$ \\
    time & $\SI{1}{s} = \SI{1.5193e15}{eV^{-1}}$ & $\SI{1}{yr} = \SI{3.156e7}{s} = \SI{4.795e22}{eV^{-1}}$ \\
    temperature & $\SI{1}{K} = \SI{8.6173e-5}{eV}$, $\SI{1}{eV} = \SI{11605}{K}$ & $\SI{300}{K} = \SI{25.85}{meV}$ \\
    speed & $\SI{1}{km.s^{-1}} = \num{3.3356e-6}$ & $v_{\mathrm{max}} = \SI{794.5}{km.s^{-1}} = \num{2.650e-3}$ \\
    number density & $\SI{1}{cm^{-3}} = \SI{7.6835e-15}{eV^3}$ & $\rho_\chi = \SI{0.3}{GeV.cm^{-3}} = \SI{2.31e-6}{eV^4}$ \\
    mass density & $\SI{1}{g.cm^{-3}} = \SI{4.310e18}{eV^4}$ & silicon, \SI{2.33}{g.cm^{-3}} $= \SI{1.004e19}{eV^4}$ \\
    mass & $\SI{1}{kg} = \SI{5.610e26}{GeV}$ & -- \\
    cross-section & $\SI{1}{cm^2} = \SI{2.5682e9}{eV^{-2}}$, $\SI{1}{b} = \SI{2.5682e-15}{eV^{-2}}$ & $\SI{1}{GeV^{-2}} = \SI{0.3894}{mb} = \SI{3.894e-28}{cm^2}$ \\
    spectroscopy & $\SI{1}{THz} = \SI{4.1357}{meV}$ ($h\nu$), $\SI{1}{cm^{-1}} = \SI{0.12398}{meV}$ & -- \\
    \bottomrule
  \end{tabular}
\end{table}

\Cref{tab:A-conversions} lists the factors. Each entry is a power of one of the products in \cref{eq:A.1-constants}. For example, $\SI{1}{cm^{-3}} = (\hbar c)^3/\si{cm^3} = (\SI{1.973270e-5}{eV})^3 = \SI{7.6835e-15}{eV^3}$, and $\SI{1}{b} = \SI{e-24}{cm^2}$ divided by $(\hbar c)^2 = \SI{3.8938e-10}{eV^2.cm^2}$ gives \SI{2.5682e-15}{eV^{-2}}. A mass density in \si{g.cm^{-3}} first becomes an energy density through $\SI{1}{g} = \SI{5.610e23}{GeV}$ and then picks up $(\hbar c)^3$. Spectroscopic frequencies need care: a frequency $\nu$ quoted in \si{THz} corresponds to the energy $h\nu = 2\pi\hbar\nu$, while an angular frequency $\omega$ in \si{rad.s^{-1}} corresponds to $\hbar\omega$. Phonon tables usually quote $\nu$, so the optical phonon of silicon at the zone center, at \SI{15.5}{THz}~\cite{nilsson1972study}, has $\omega_{Q\nu} = \SI{64.1}{meV}$.

\subsection{Checking dimensions}

In natural units every quantity has a single dimension, a power of energy, and a formula is checked by counting powers of eV. The rate per unit target mass of \cref{eq:1.5-rate} is the standard example: it is dimensionless, as a rate (eV) per unit mass (eV) must be (\cref{sec:1.5}). The dimensions of the main objects are fixed by their definitions: the dynamic structure factor $S(q,\omega)$ has dimension $\mathrm{eV}^2$ (\cref{sec:4.2}), the transform $\tilde V(\vect{q})$ of a potential has $\mathrm{eV}^{-2}$, the density vertex $\rho_{mn}(\vect{k},\vect{q})$ and the dielectric function are dimensionless, and $\sigmae$ has $\mathrm{eV}^{-2}$. Final numbers are always converted back: rates to \si{kg^{-1}.yr^{-1}}, cross-sections to \si{cm^2}, momenta to \si{keV}, and energies to \si{eV} or \si{meV}. A rate per unit mass in natural units is converted with $\SI{1}{kg^{-1}.yr^{-1}} = \num{3.718e-59}$, the inverse of the product of \SI{1}{kg} in eV and \SI{1}{yr} in $\mathrm{eV}^{-1}$.

The metric signature is $(+,-,-,-)$, with $x^\mu = (t, \vect{x})$ and $\partial_\mu\partial^\mu = \partial_t^2 - \nabla^2$, so that a free particle obeys $p^\mu p_\mu = E^2 - |\vect{p}|^2 = m^2$; it is the signature of most particle-physics texts, and \cref{app:C.9} uses it. Physical values in the text are quoted in conventional units with a thin space, \SI{1.17}{eV} or \SI{0.3}{GeV.cm^{-3}}, and natural-unit values are written without units only for speeds.
